\section{Proofs for Spectral Targets and Structured Families}
\label{app:matrix-proofs}

\begin{proof}[Proof of \cref{thm:full-spd-benchmark}]
Let \(S_1\in\Ck\), and let \(z_1\le\cdots\le z_d\) be the ordered
log-eigenvalues of \(S_1\). The affine-invariant distance satisfies the
spectral lower bound
\[
  d_{\AI}(S_0,S_1)\ge \norm{y-z}_2,
\]
where \(y\) is the ordered log-spectrum of \(S_0\). This is the standard
eigenvalue majorization lower bound for the affine-invariant metric on
positive-definite matrices \cite{bhatia2007positive}. Since \(S_1\in\Ck\), the
vector \(z\) has width at most \(\log K\), so \(z_i\in[c,c+\log K]\) for some
\(c\). Hence
\[
  d_{\AI}(S_0,S_1)^2
  \ge
  \sum_{i=1}^d \dist(y_i,[c,c+\log K])^2.
\]
Taking the infimum over \(S_1\in\Ck\) gives the lower bound.

For the reverse inequality, fix \(c\) and define
\[
  z_i(c)=\Pi_{[c,c+\log K]}(y_i),
\]
the scalar projection of \(y_i\) onto the interval. The vector \(z(c)\) has
width at most \(\log K\). Choose \(S_1\) with the same eigenvectors as \(S_0\)
and log-eigenvalues \(z_i(c)\). Then \(S_1\in\Ck\) and the commuting
affine-invariant distance is exactly \(\norm{y-z(c)}_2\). Minimizing over
\(c\in\R\) proves the formula.
\end{proof}

\begin{proof}[Proof of \cref{prop:monotonicity}]
Under the stated compatibility assumption, every admissible \(\F_1\)-path from
the common initial point \(S_0\) to the \(\F_1\)-target is also an admissible
\(\F_2\)-path with the same length. Moreover
\(\mathcal C_{K,\F_1}(H)\subseteq\mathcal C_{K,\F_2}(H)\). Thus the feasible
path set for \(\F_2\) contains the feasible path set for \(\F_1\), and the
infimum cannot increase. Enlarging the target threshold \(K\) similarly
enlarges the feasible endpoint set. The ambient SPD cone contains every
admissible restricted path considered in the lower-bound statement, giving
\(D_{K,\F}\ge D_K\).
\end{proof}

\begin{proof}[Proof of \cref{prop:submanifold-distance}]
Under the smooth embedded submanifold assumption, admissible paths are exactly
absolutely continuous paths in \(\mathcal S_\F(H)\) from \(S_0\) to
\(\mathcal C_{K,\F}(H)\). The induced length is the ambient affine-invariant
length restricted to that submanifold. Hence the infimum is the intrinsic
distance to the target. If no path connects the component of \(S_0\) to the
target, the intrinsic distance is \(+\infty\).
\end{proof}

\begin{proof}[Proof of \cref{thm:scale-closed-threshold}]
We first record the sharp vector inequality used below. For
\(x\in\R^d\), write
\[
  \operatorname{osc}(x)=\max_i x_i-\min_i x_i,\qquad
  \dist(x,\R\mathbf 1)=\min_{s\in\R}\norm{x-s\mathbf 1}_2.
\]
Let \(\tau=\operatorname{osc}(x)\). Translating \(x\) by a scalar multiple of
\(\mathbf 1\) changes neither side, so assume \(\min_i x_i=0\) and
\(\max_i x_i=\tau\). Then \(0\le x_i\le\tau\). The squared distance to
\(\R\mathbf 1\) is \(\sum_i(x_i-\bar x)^2\), which is the variance numerator.
Among vectors in the box \([0,\tau]^d\), this convex function is maximized at a
vertex. If \(k\) coordinates equal \(\tau\) and \(d-k\) equal \(0\), then
\[
  \sum_i(x_i-\bar x)^2=\frac{k(d-k)}{d}\tau^2
  \le\frac{\lfloor d^2/4\rfloor}{d}\tau^2.
\]
Thus
\[
  \operatorname{osc}(x)
  \ge
  \sqrt{\frac{d}{\lfloor d^2/4\rfloor}}
  \dist(x,\R\mathbf 1).
\]
The constant is attained by placing \(\lfloor d/2\rfloor\) coordinates at one
endpoint and the remaining coordinates at the other endpoint.

Now fix \(Q\in\F\) and let
\[
  \lambda_i(Q^{-1/2}HQ^{-1/2})=\exp(\ell_i).
\]
Then
\[
  \log\kappa(Q^{-1}H)=\operatorname{osc}(\ell).
\]
Because \(\F\) is scale closed, \(cQ\in\F\) for every \(c>0\). The eigenvalues
of \(H^{-1/2}(cQ)H^{-1/2}\) are \(\exp(\log c-\ell_i)\), so
\[
  \inf_{c>0}d_{\AI}(H,cQ)
  =
  \dist(\ell,\R\mathbf 1).
\]
The vector inequality gives
\[
  \log\kappa(Q^{-1}H)
  \ge
  \alpha_d\inf_{c>0}d_{\AI}(H,cQ)
  \ge
  \alpha_d d_{\AI}(H,\F).
\]
Taking the infimum over \(Q\in\F\) proves the lower bound. Sharpness follows
from the scale family \(\F=\{cI:c>0\}\) and any equality vector in the sharp
oscillation-distance inequality, used as the log-spectrum of \(H\).
\end{proof}

\begin{proof}[Proof of
\cref{thm:structured-log-spectral-mother,thm:restricted-geometrodynamic-mother}]
We first prove the log-majorization statement that drives the theorem. For
\(A,B\in\SPD^d\), \(t\in[0,1]\), and \(1\le k\le d\), the exterior-power
representation commutes with the weighted geometric mean:
\[
  \bigwedge^k(A\#_tB)
  =
  \left(\bigwedge^kA\right)\#_t
  \left(\bigwedge^kB\right).
\]
For positive-definite \(X,Y\), monotonicity of the geometric mean gives
\[
  \norm{X\#_tY}_{\mathrm{op}}
  \le
  \norm{X}_{\mathrm{op}}^{1-t}
  \norm{Y}_{\mathrm{op}}^t,
\]
because \(X\preceq\norm{X}_{\mathrm{op}}I\) and
\(Y\preceq\norm{Y}_{\mathrm{op}}I\). Applying this inequality to the
exterior powers yields
\[
  \prod_{i=1}^k\lambda_i^\downarrow(A\#_tB)
  \le
  \prod_{i=1}^k
  \lambda_i^\downarrow(A)^{1-t}
  \lambda_i^\downarrow(B)^t.
\]
Equality holds for \(k=d\) because
\(\det(A\#_tB)=\det(A)^{1-t}\det(B)^t\). Taking logarithms proves
\begin{equation}
  \log\lambda(A\#_tB)
  \prec
  (1-t)\log\lambda(A)+t\log\lambda(B),
  \tag{A.1}
\end{equation}
up to the irrelevant common choice of decreasing order.

Fix an index \(a\), abbreviate \(H=H_a\) and \(\phi=\phi_a\), and set
\[
  s_H(G)=\log\lambda(G^{-1/2}HG^{-1/2}).
\]
Now fix
\(G_0,G_1\in\F\) and let \(G_t=G_0\#_tG_1\). Congruence by
\(H^{-1/2}\) is an affine-invariant isometry and preserves geometric means, so
with \(A_i=H^{-1/2}G_iH^{-1/2}\) we have
\[
  H^{-1/2}G_tH^{-1/2}=A_0\#_tA_1.
\]
The eigenvalues of \(G^{-1/2}HG^{-1/2}\) are the reciprocals of those of
\(H^{-1/2}GH^{-1/2}\). Majorization is preserved by multiplication by
\(-1\), and a convex permutation-invariant function is Schur convex.
Therefore (A.1) and ordinary convexity of \(\phi\) give
\[
\begin{aligned}
  \Omega_{\phi,H}(G_t)
  &\le
  \phi\!\left((1-t)s_H(G_0)+t s_H(G_1)\right)\\
  &\le
  (1-t)\Omega_{\phi,H}(G_0)+t\Omega_{\phi,H}(G_1).
\end{aligned}
\]
This proves geodesic convexity of each component \(\Omega_a\). Since
\(\rho\) is coordinatewise nondecreasing,
\[
  \rho(\Omega_1(G_t),\ldots,\Omega_q(G_t))
  \le
  \rho((1-t)\Omega(G_0)+t\Omega(G_1)),
\]
and convexity of \(\rho\) proves geodesic convexity of \(\mathcal J\).
All functions in the statement are finite convex functions on Euclidean
spaces and hence continuous, so \(\mathcal J\) is closed.

If \(\phi_a(x+c\mathbf1)=\phi_a(x)\), then
\[
  \log\lambda((cG)^{-1/2}H_a(cG)^{-1/2})
  =s_{H_a}(G)-(\log c)\mathbf1,
\]
which proves scale invariance. For stability, the standard affine-invariant
spectral inequality gives
\[
  \norm{s_{H_a}(G)-s_{\widehat H_a}(G)}_2
  \le
  d_{\AI}(H_a,\widehat H_a),
\]
because congruence by \(G^{-1/2}\) is an isometry. The Lipschitz bounds for
\(\phi_a\) and \(\rho\), followed by Cauchy--Schwarz, give the displayed
uniform perturbation estimate.

  We now prove the static projection theorem.  For quadratic kinetic action,
  \(U=0\) and
\(\Gamma=\iota_{\mathcal C_\tau}\), every curve over the
remaining horizon \(h=T-t\) satisfies
\[
  \frac12\int_t^T\norm{\dot\gamma}^2\dd s
  \ge\frac{\Len(\gamma)^2}{2h}
  \ge\frac{d(G,\mathcal C_\tau)^2}{2h}.
\]
Equality is attained by the constant-speed geodesic to the metric projection,
which proves the closed form for \(V\).

A sublevel set of a closed geodesically convex function on a closed
geodesically convex family is closed
and geodesically convex. In a finite-dimensional Hadamard manifold it has a
unique metric projection: existence follows by restricting a minimizing
sequence to a closed bounded ball and applying Hopf--Rinow, while uniqueness
follows from strict geodesic convexity of squared distance. The geodesic from
\(G_0\) to the projection stays in \(\F\) and reaches the target with length
equal to the metric distance. Every admissible curve has length at least the
distance between its endpoints, proving the least-action identity and the
uniqueness of the constant-speed minimizer. The displayed Pythagorean
inequality is the standard projection inequality in a Hadamard manifold.
Family and threshold monotonicity follow from inclusion of the corresponding
target and path sets. Finally,
\[
  |\widehat{\mathcal J}-\mathcal J|\le\varepsilon
\]
immediately gives the two target inclusions; applying target monotonicity gives
the perturbation bracket for the dynamic cost.

It remains to prove exactness and convergence of the algorithm. Let
\(P=P_\tau\). Since \(G_s\) is feasible,
\(d(G_0,P)\le R\), so \(P\in\mathcal D\). Closed balls in a Hadamard
manifold are geodesically convex, hence \(\mathcal D\) is closed and
geodesically convex. Write
\[
  f(G)=\frac12d(G_0,G)^2,
  \qquad h(G)=\mathcal J(G)-\tau.
\]
Take any infeasible \(G\in\mathcal D\), set \(a=h(G)>0\), and follow the
geodesic from \(G\) to \(G_s\). At
\[
  t=\frac{a}{a+\sigma}
\]
the point \(Y=G\#_tG_s\) is feasible, because geodesic convexity gives
\(h(Y)\le(1-t)a-t\sigma=0\). Both endpoints lie in the radius-\(R\) ball, so
\[
  d(G,Y)
  =t\,d(G,G_s)
  \le\frac{2Ra}{a+\sigma}
  \le\frac{2Ra}{\sigma}.
\]
On this ball the function \(f\) is \(R\)-Lipschitz, since
\[
  |f(X)-f(Y)|
  \le
  \frac{d(G_0,X)+d(G_0,Y)}{2}d(X,Y)
  \le R d(X,Y).
\]
Consequently
\[
  \Phi_\Lambda(G)
  =f(G)+\Lambda a
  \ge
  f(Y)+\left(\Lambda-\frac{2R^2}{\sigma}\right)a
  >f(Y)\ge f(P)=\Phi_\Lambda(P).
\]
No infeasible point minimizes \(\Phi_\Lambda\). On the feasible set the
penalty vanishes and \(P\) uniquely minimizes \(f\), proving exactness.

The function \(f\) is \(1\)-strongly geodesically convex on a Hadamard
manifold, while \(\pospart{h}=\max\{h,0\}\) is geodesically convex. Thus
\(\Phi_\Lambda\) is \(1\)-strongly geodesically convex on \(\mathcal D\).
Each proximal objective is proper, closed, and
\((1+\eta^{-1})\)-strongly geodesically convex, so it has a unique minimizer
\(G_{r+1}\in\mathcal D\). Strong convexity at the two minimizers
\(P=\argminop_{\mathcal D}\Phi_\Lambda\) and \(G_{r+1}\) gives
\[
\begin{aligned}
  &\Phi_\Lambda(P)+\frac{1}{2\eta}d(G_r,P)^2\\
  &\quad\ge
  \Phi_\Lambda(G_{r+1})
  +\frac{1}{2\eta}d(G_r,G_{r+1})^2
  +\frac{1+\eta^{-1}}{2}d(G_{r+1},P)^2\\
  &\quad\ge
  \Phi_\Lambda(P)
  +\frac12d(G_{r+1},P)^2
  +\frac{1}{2\eta}d(G_r,G_{r+1})^2
  +\frac{1+\eta^{-1}}{2}d(G_{r+1},P)^2.
\end{aligned}
\]
After cancellation and multiplication by \(2\eta\),
\[
  d(G_r,P)^2
  \ge
  (1+2\eta)d(G_{r+1},P)^2+d(G_r,G_{r+1})^2.
\]
Dropping the last term and iterating proves the stated global contraction.
\end{proof}

\begin{proof}[Proof of \cref{lem:log-spectral-pullback}]
Put \(Y=\log S\) and
\[
  S_t=e^{-tX/2}Se^{-tX/2},
  \qquad
  \dot S_0=-\frac12(XS+SX).
\]
The matrix logarithm is Fr\'echet differentiable on the SPD cone, so
\[
  \frac{\dd}{\dd t}\bigg|_{t=0}\log S_t
  =D\log_S[\dot S_0].
\]
For the finite convex spectral function \(F\), convex directional calculus
gives
\[
  F'(Y;D\log_S[\dot S_0])
  =
  \max_{Z\in\partial F(Y)}
  \tr\!\left(ZD\log_S[\dot S_0]\right).
\]
The spectral subdifferential theorem \cite{lewis1996convex} implies that every
\(Z\in\partial F(Y)\) commutes with \(Y\), hence also with \(S=e^Y\).
The Fr\'echet derivative \(D\log_S\) is self-adjoint in the Frobenius inner
product, and commutation gives
\[
  D\log_S[Z]=S^{-1}Z.
\]
Therefore
\[
\begin{aligned}
  \tr\!\left(ZD\log_S[\dot S_0]\right)
  &=\tr\!\left(S^{-1}Z\dot S_0\right)\\
  &=-\frac12\tr\!\left(S^{-1}ZXS+S^{-1}ZSX\right)\\
  &=-\tr(ZX).
\end{aligned}
\]
Substitution into the support-function formula proves the directional identity.
Since a closed convex subdifferential is determined by its support function,
the whitened AIRM subdifferential is exactly \(-\partial F(Y)\), with no
simplicity assumption on the spectrum.
\end{proof}

\begin{proof}[Proof of \cref{thm:normal-response-law}]
Represent a tangent direction at \(G\) in whitened coordinates by
\(X\in\widehat T_G\F\), so the corresponding geodesic is
\(G_t=G^{1/2}\exp(tX)G^{1/2}\).  The generalized eigenvalues of
\((H_a,G_t)\) are the eigenvalues of
\[
  \exp(-tX/2)S_a\exp(-tX/2).
\]
Applying \cref{lem:log-spectral-pullback} gives, at simple or repeated spectra,
\[
  \partial\Omega_a(G)=-\mathfrak Z_a(G)
\]
in whitened AIRM coordinates.

The finite convex composition rule for the coordinatewise nondecreasing
aggregator \(\rho\) now yields
\[
  \partial\mathcal J(G)
  =
  \left\{-\sum_{a=1}^q\theta_aZ_a:
  \theta\in\partial\rho(\Omega(G)),\
  Z_a\in\mathfrak Z_a(G)\right\}
\]
in whitened coordinates.  On the assumed closed geodesically convex smooth
submanifold, the geodesic subgradient inequality shows that an ambient
subgradient orthogonal to the tangent space is globally sufficient.  The
constrained Fermat rule gives the converse.  Since a normal space is a linear
space, the minus sign is immaterial, proving the aggregate normal-response
criterion.

For \(P=P_\F(H)\), the whitened gradient at \(P\) of the squared-distance
objective is
\[
  -\log(P^{-1/2}HP^{-1/2}).
\]
Projection optimality therefore gives
\(\log S\in\widehat N_P\F\). Applying the first part with \(q=1\) proves the
nearest-versus-best coincidence criterion.
\end{proof}

\begin{proof}[Proof of \cref{lem:diag-distance-identity}]
Congruence by \(H^{-1/2}\) is an isometry for \(d_{\AI}\), so it is enough to
compute distances between diagonal matrices. For
\(D(x)=\diag(e^{x_1},\ldots,e^{x_d})\),
\[
  d_{\AI}(D(x),D(y))
  =
  \fro{\log(D(x)^{-1/2}D(y)D(x)^{-1/2})}
  =
  \norm{x-y}_2,
\]
because the matrices commute. The same calculation applied to absolutely
continuous paths gives the Euclidean line element in \(x\)-coordinates. The
target \(\mathcal X_K(H)\) is the preimage of the closed set
\(\Ck\subset\SPD^d\) under the continuous map
\(x\mapsto H^{-1/2}D(x)H^{-1/2}\), hence is closed. Therefore the restricted
complexity is exactly the Euclidean distance from \(x_0\) to this closed target
set, with the usual convention that the distance to the empty set is
\(+\infty\).
\end{proof}

\begin{proof}[Proof of \cref{thm:diag-lmi}]
Suppose \(\kappa(D^{-1}H)\le K\). Let
\[
  \lambda_{\min}=\min_{v\ne0}\frac{v^\top Hv}{v^\top Dv}.
\]
Then all generalized Rayleigh quotients lie in
\([\lambda_{\min},K\lambda_{\min}]\). With \(E=\lambda_{\min}D\),
\[
  E\preceq H\preceq K E.
\]
Conversely, if \(E\preceq H\preceq KE\) for a positive diagonal \(E\), then
every generalized Rayleigh quotient \(v^\top Hv/v^\top Ev\) lies in
\([1,K]\), so \(\kappa(E^{-1}H)\le K\).
\end{proof}

\begin{proof}[Proof of \cref{cor:diag-aligned}]
If \(H\) and \(G_0\) are diagonal, then \(S_0=H^{-1/2}G_0H^{-1/2}\) is diagonal.
The full SPD projection onto \(\Ck\) keeps the eigenvectors of \(S_0\) and
clips only its log-spectrum. Therefore the optimal endpoint and the
affine-invariant geodesic from \(S_0\) to that endpoint remain diagonal. The
diagonal family contains a full SPD shortest path, and \cref{prop:monotonicity}
gives equality.
\end{proof}

\begin{proof}[Proof of \cref{thm:block-lmi}]
The proof is identical to \cref{thm:diag-lmi}, replacing the diagonal cone by
the block-diagonal cone. If \(\kappa(G^{-1}H)\le K\), take
\(E=\lambda_{\min}G\), which is block diagonal. The converse follows from the
generalized Rayleigh quotient bound.
\end{proof}

\begin{proof}[Proof of \cref{prop:structured-dual-certificate}]
Assume, for contradiction, that there exists \(E\in\mathcal L\) with
\[
  H-E\succeq0,\qquad KE-H\succeq0.
\]
For \(P,Q\succeq0\),
\[
  0
  \le
  \ip{P}{H-E}+\ip{Q}{KE-H}
  =
  \ip{P-Q}{H}+\ip{-P+KQ}{E}.
\]
Since \(E\in\mathcal L\) and \(\Pi_{\mathcal L}(P-KQ)=0\), the last term is
zero. Thus any feasible \(E\) would imply \(\ip{P-Q}{H}\ge0\), contradicting
the strict inequality in the certificate.
\end{proof}

\begin{proof}[Proof of \cref{cor:sign-flip-certificate}]
Let \(q\) be a unit eigenvector with
\[
  q^\top(KTHT-H)q<0.
\]
Set
\[
  Q=qq^\top,\qquad P=K(Tq)(Tq)^\top=KTQT.
\]
Then \(P,Q\succeq0\). Because \(T\) is a constant sign on each block,
\(\Pi_{\mathcal L}(TQT-Q)=0\), and hence
\[
  \Pi_{\mathcal L}(P-KQ)=K\Pi_{\mathcal L}(TQT-Q)=0.
\]
Moreover,
\[
  \ip{P-Q}{H}
  =
  q^\top(KTHT-H)q
  <0.
\]
\Cref{prop:structured-dual-certificate} proves infeasibility.
\end{proof}

\begin{proof}[Proof of \cref{prop:block-jacobi}]
The block Jacobi matrix \(G_{\mathrm{BJ}}\) is feasible in the block family, so
the best block condition number is at most the one it achieves. If
\(\opnorm{\widetilde H-I}\le\delta<1\), then
\[
  (1-\delta)I\preceq \widetilde H\preceq (1+\delta)I,
\]
which gives the stated condition-number bound.
\end{proof}

\begin{proof}[Proof of \cref{prop:kron-spec-lower}]
The restricted target
\[
  (Y+\mathcal A_{\mathrm{kron}})\cap\Yk
\]
is a subset of \(\Yk\). Distance to a subset is at least distance to the full
set. If the full projection lies in the restricted target, it attains both
distances.
\end{proof}

\begin{proof}[Proof of \cref{prop:kron-line-element}]
For \(G=B\otimes A\),
\[
  G^{-1/2}\dot G\,G^{-1/2}
  =
  Y\otimes I_m+I_n\otimes X.
\]
Its Frobenius norm squared is
\[
  m\fro{Y}^2+n\fro{X}^2+2\tr(X)\tr(Y),
\]
which is the displayed formula. For \(X=\alpha I_m\) and \(Y=-\alpha I_n\), the
three terms sum to zero.
\end{proof}

\begin{proof}[Proof of \cref{prop:kron-product-distance}]
Using Kronecker identities,
\[
  G_0^{-1/2}G_1G_0^{-1/2}
  =
  (B_0^{-1/2}B_1B_0^{-1/2})
  \otimes
  (A_0^{-1/2}A_1A_0^{-1/2}).
\]
Taking the logarithm gives
\[
  \log(G_0^{-1/2}G_1G_0^{-1/2})
  =
  L_B\otimes I_m+I_n\otimes L_A.
\]
The squared Frobenius norm of this matrix is
\[
  m\fro{L_B}^2+n\fro{L_A}^2+2\tr(L_A)\tr(L_B),
\]
which is the claimed formula.
\end{proof}

\begin{proof}[Proof of \cref{prop:kron-fixed-basis}]
In the fixed Kronecker eigenbasis, the eigenvalues of \(X_t\) and \(Y_t\) are
\(\dot\alpha_i\) and \(\dot\beta_j\). By \cref{prop:kron-line-element},
\[
  \norm{\dot G_t}_{G_t,\AI}^2
  =
  n\sum_i\dot\alpha_i^2+
  m\sum_j\dot\beta_j^2+
  2\left(\sum_i\dot\alpha_i\right)\left(\sum_j\dot\beta_j\right).
\]
On the other hand,
\[
  \fro{\dot Z_t}^2
  =
  \sum_{i,j}(\dot\alpha_i+\dot\beta_j)^2,
\]
which expands to the same expression.
\end{proof}

\begin{proof}[Proof of \cref{thm:kron-log-projection}]
If \(M=B\otimes A\), then the spectral calculus for Kronecker products gives
\[
  \log M=(\log B)\otimes I_m+I_n\otimes(\log A)\in\mathcal L_{m,n}.
\]
Conversely, if
\[
  \log M=Y\otimes I_m+I_n\otimes X,
\]
then the two summands commute, and therefore
\[
  M=\exp(Y\otimes I_m)\exp(I_n\otimes X)=\exp(Y)\otimes\exp(X).
\]
This proves the log characterization.

For the projection formula, decompose
\[
  \mathcal L_{m,n}
  =
  \operatorname{span}\{I_{mn}\}
  \oplus
  \{I_n\otimes X:\tr X=0\}
  \oplus
  \{Y\otimes I_m:\tr Y=0\},
\]
an orthogonal direct sum under the Frobenius inner product. With
\[
  P=\tau I_{mn}+I_n\otimes X_0+Y_0\otimes I_m,
\]
the definitions of \(\tau,X_0,Y_0\) give
\[
  \tr P=\tr L,\qquad
  \operatorname{Tr}_B(P)=\operatorname{Tr}_B(L),\qquad
  \operatorname{Tr}_A(P)=\operatorname{Tr}_A(L).
\]
Thus \(L-P\) is orthogonal to each component of \(\mathcal L_{m,n}\), so
\(P=\Pi_{\mathcal L}L\). The residual statement follows by applying the log
characterization to \(L=\log M\).
\end{proof}

\begin{proof}[Proof of \cref{lem:hadamard-projection-facts}]
Closed geodesically convex subsets of a Hadamard manifold are proximinal and
the projection is unique because squared distance is strictly convex along
geodesics. The same CAT(0) convexity inequality gives, for any geodesic
\(\gamma\) in \(C\),
\[
  f(\gamma_t)
  \le
  (1-t)f(\gamma_0)+tf(\gamma_1)
  -\frac12t(1-t)d(\gamma_0,\gamma_1)^2,
\]
which is \(1\)-strong geodesic convexity of \(f(y)=d(y,x)^2/2\).

For the line-search statement, the descent curve is a geodesic initialized in
the negative restricted-gradient direction. On a compact geodesic neighborhood
where the restricted gradient is \(L\)-Lipschitz, the geodesic descent lemma
gives
\[
  f(\gamma_\eta)
  \le
  f(\gamma_0)-\eta\norm{\operatorname{grad}_C f(\gamma_0)}^2
  +\frac{L}{2}\eta^2\norm{\operatorname{grad}_C f(\gamma_0)}^2.
\]
Thus all sufficiently small positive \(\eta\) satisfy the Armijo inequality.
Compactness makes the admissible upper step size uniform over the sublevel set,
and geometric backtracking then accepts a step bounded below by a positive
constant depending only on the Lipschitz constant and the backtracking
parameters.
\end{proof}

\begin{proof}[Proof of \cref{thm:kron-ai-projection}]
By \cref{thm:kron-log-projection},
\(\Kron_{m,n}=\exp(\mathcal L_{m,n})\). The principal matrix logarithm is a
global diffeomorphism from \(\SPD^{mn}\) to \(\mathbb S^{mn}\), and
\(\mathcal L_{m,n}\) is a linear subspace. Hence \(\Kron_{m,n}\) is a smooth
embedded submanifold. This formulation also removes the factor-scale gauge:
different pairs \((X,Y)\) that differ by
\((X+\alpha I_m,Y-\alpha I_n)\) represent the same element of
\(\mathcal L_{m,n}\).

For \(G_i=B_i\otimes A_i\), the affine-invariant geodesic satisfies
\[
  G_0^{-1/2}G_1G_0^{-1/2}
  =
  (B_0^{-1/2}B_1B_0^{-1/2})
  \otimes
  (A_0^{-1/2}A_1A_0^{-1/2}).
\]
Since \(\exp(t\log(P\otimes Q))=P^t\otimes Q^t\) for \(P,Q\succ0\), the
ambient geodesic between \(G_0\) and \(G_1\) remains in \(\Kron_{m,n}\). Hence
\(\Kron_{m,n}\) is totally geodesic.

It is also closed. If \(B_k\otimes A_k\to G\succ0\), fix the gauge
\(\det A_k=1\). Uniform lower and upper eigenvalue bounds on
\(B_k\otimes A_k\) bound all eigenvalue ratios of \(A_k\); the determinant
gauge then bounds the eigenvalues of \(A_k\) above and below, and the product
eigenvalue bounds do the same for \(B_k\). Passing to a subsequence gives
\(A_k\to A\succ0\) and \(B_k\to B\succ0\), hence \(G=B\otimes A\).

The SPD cone with \(d_{\AI}\) is a Hadamard manifold. By
\cref{lem:hadamard-projection-facts}, projection onto a closed geodesically
convex subset is unique, giving \(G_\star\).
For \(G=B\otimes A\), tangent vectors in whitened coordinates are exactly
\[
  Y\otimes I_m+I_n\otimes X\in\mathcal L_{m,n}.
\]
Indeed, every such \(Z\in\mathcal L_{m,n}\) generates the curve
\(G^{1/2}\exp(tZ)G^{1/2}\in\Kron_{m,n}\), while differentiating any smooth
factor curve \(B_t\otimes A_t\) and whitening by \(G^{-1/2}\) gives an element
of this same subspace. Thus the normal space is the Frobenius orthogonal
complement of \(\mathcal L_{m,n}\) in whitened coordinates.
At the projection point, the logarithmic residual
\[
  R_\star=\log(G_\star^{-1/2}MG_\star^{-1/2})
\]
is orthogonal to every tangent direction. Orthogonality to
\(I_n\otimes X\) and \(Y\otimes I_m\) is exactly
\[
  \operatorname{Tr}_B(R_\star)=0,\qquad
  \operatorname{Tr}_A(R_\star)=0.
\]
Conversely, these equations give the first-order projection condition on a
closed geodesically convex set in a Hadamard manifold, hence the unique
projection. The distance formula is the definition of \(d_{\AI}\).
\end{proof}

\begin{proof}[Proof of \cref{cor:kron-residual-certificate}]
The equivalence \(V(G)=0\Longleftrightarrow G=P_{\Kron}(M)\) is the normal
equation in \cref{thm:kron-ai-projection}. The quantitative bounds use
\cref{lem:hadamard-projection-facts}: the objective
\[
  f(G)=\frac12d_{\AI}(G,M)^2
\]
is \(1\)-strongly geodesically convex on the totally geodesic Hadamard
submanifold \(\Kron_{m,n}\). Let \(G_\star=P_{\Kron}(M)\),
\(v=\operatorname{Log}_G(G_\star)\), and
\(d=\norm{v}=d_{\AI}(G,G_\star)\). Strong convexity along the geodesic from
\(G\) to \(G_\star\) gives
\[
  f(G)-f(G_\star)\le \fro{V(G)}d-\frac12d^2.
\]
Strong convexity at the minimizer gives
\[
  f(G)-f(G_\star)\ge\frac12d^2.
\]
Combining the inequalities yields \(d\le\fro{V(G)}\), and maximizing
\(\fro{V(G)}d-d^2/2\) over \(0\le d\le\fro{V(G)}\) gives the objective-gap
bound.
\end{proof}

\begin{proof}[Proof of \cref{cor:kron-armijo-solver}]
If \(G_r=B_r\otimes A_r\), then \(V_r\in\mathcal L_{m,n}\) has the form
\[
  V_r=Y_r\otimes I_m+I_n\otimes X_r.
\]
The two summands commute, so \(\exp(\eta V_r)=\exp(\eta Y_r)\otimes
\exp(\eta X_r)\), and the trial point remains in \(\Kron_{m,n}\).

Let \(G_\star=P_{\Kron}(M)\). The restricted negative gradient of
\[
  f(G)=\frac12d_{\AI}(G,M)^2
\]
is represented in whitened coordinates by
\(V(G)=\Pi_{\mathcal L}R_G(M)\). The initial sublevel set
\(\mathcal Q_0=\{G\in\Kron_{m,n}:f(G)\le f(G_0)\}\) is closed and bounded in
the complete finite-dimensional Hadamard submanifold \(\Kron_{m,n}\), hence
compact by Hopf--Rinow. Since the SPD cone has no cut locus and the
squared-distance objective is smooth, the restricted gradient is Lipschitz on
a compact geodesic neighborhood of \(\mathcal Q_0\); let \(L_0\) be such a
constant. By \cref{lem:hadamard-projection-facts}, Armijo backtracking
terminates and the accepted step sizes have a positive lower bound
\(\eta_{\min}>0\) depending only on
\(\mathcal Q_0,L_0,\bar\eta,\gamma\), and \(c\)
\cite{absil2008optimization,boumal2023introduction}.

Armijo decrease gives
\[
  f(G_{r+1})\le f(G_r)-c\eta_r\fro{V_r}^2.
\]
Thus \(f(G_r)\) decreases and \(\sum_r\eta_r\fro{V_r}^2<\infty\). Since
\(\eta_r\ge\eta_{\min}\), we have \(\fro{V_r}\to0\). Any cluster point
\(\bar G\) in the compact sublevel set satisfies \(V(\bar G)=0\), hence
\(\bar G=G_\star\) by \cref{cor:kron-residual-certificate}. The cluster point is
unique, so the full sequence converges to \(G_\star\).

For the rate, set
\[
  \underline\eta=\min\{\eta_{\min},(2c)^{-1}\}.
\]
\Cref{cor:kron-residual-certificate} gives
\[
  \fro{V_r}^2\ge2\bigl[f(G_r)-f(G_\star)\bigr].
\]
Combining this with Armijo decrease and \(\eta_r\ge\underline\eta\) yields
\[
  f(G_{r+1})-f(G_\star)
  \le
  (1-2c\underline\eta)\bigl[f(G_r)-f(G_\star)\bigr],
\]
and iteration proves the displayed bound.
\end{proof}

\begin{proof}[Proof of \cref{thm:kron-k-target}]
For the self-conditioned metric target, the condition-number set \(\Ck\) is
closed and geodesically convex because
the weighted matrix geometric mean satisfies
\[
  \lambda_{\max}(S_0\#_tS_1)
  \le
  \lambda_{\max}(S_0)^{1-t}\lambda_{\max}(S_1)^t
\]
and the analogous lower bound for \(\lambda_{\min}\). Together with
\cref{thm:kron-ai-projection}, this makes
\(\Kron_{m,n,K}=\Kron_{m,n}\cap\Ck\) closed and geodesically convex.

The two displayed distance bounds are the Hadamard Pythagorean inequality for
projection onto \(\Kron_{m,n}\), followed by the triangle inequality through
\(G_\star=P_{\Kron}(M)\).

It remains to prove the QP formula. For \(B\otimes A\),
\[
  \kappa(B\otimes A)=\kappa(B)\kappa(A),
\]
so the target condition is
\[
  (x_m-x_1)+(y_n-y_1)\le\log K
\]
for sorted log-eigenvalues \(x_i\) of \(A\) and \(y_j\) of \(B\). By
\cref{prop:kron-product-distance}, the squared distance between Kronecker products
is
\[
  n\,d_{\AI}(A_0,A)^2+m\,d_{\AI}(B_0,B)^2
  +2(\log\det A-\log\det A_0)(\log\det B-\log\det B_0).
\]
The eigenvalue lower bound for the affine-invariant metric and
Hoffman--Wielandt give
\[
  d_{\AI}(A_0,A)^2\ge\sum_i(x_i-a_i)^2,\qquad
  d_{\AI}(B_0,B)^2\ge\sum_j(y_j-b_j)^2.
\]
The determinant terms are the corresponding sums of log-eigenvalue
differences. Expanding gives the lower bound
\[
  \sum_{i,j}(x_i+y_j-a_i-b_j)^2.
\]
Equality is attained by choosing \(A\) and \(B\) with the same eigenvectors as
\(A_0\) and \(B_0\), proving the convex QP.
\end{proof}

\begin{proof}[Proof of \cref{thm:kron-relative-reachability}]
For any \(G\succ0\),
\[
  \kappa(G^{-1}H)
  =
  \kappa(H^{-1/2}GH^{-1/2}),
\]
because \(G^{-1}H\) is similar to
\(H^{1/2}G^{-1}H^{1/2}=(H^{-1/2}GH^{-1/2})^{-1}\), and a positive-definite
matrix and its inverse have the same spectral condition number. Thus
\(G\in\mathcal R_{K,\Kron}(H)\) if and only if
\(H^{-1/2}GH^{-1/2}\in\mathcal C_{K,\Kron}(H)\), proving the equivalence of
the first two items.

For fixed \(G\succ0\), the condition \(\kappa(G^{-1}H)\le K\) is equivalent to
the existence of \(\lambda>0\) such that
\[
  \lambda G\preceq H\preceq K\lambda G.
\]
Indeed, take
\[
  \lambda=\lambda_{\min}(G^{-1/2}HG^{-1/2})
\]
for the forward implication, and use generalized Rayleigh quotients for the
reverse implication. Setting \(G=B\otimes A\) gives the Kronecker Loewner
sandwich. The scalar \(\lambda\) can be absorbed into \(A\) or \(B\), giving
the scale-free form.

It remains to identify the distance. The target can be written as
\[
  \mathcal R_{K,\Kron}(H)
  =
  \Kron_{m,n}\cap H^{1/2}\Ck H^{1/2}.
\]
The set \(\Ck\) is closed and geodesically convex, and congruence by \(H^{1/2}\)
is an affine-invariant isometry, so \(H^{1/2}\Ck H^{1/2}\) is also closed and
geodesically convex. Intersecting with the closed totally geodesic submanifold
\(\Kron_{m,n}\) gives a closed geodesically convex subset of the Kronecker
manifold. The Kronecker manifold is a complete Hadamard submanifold under the
induced affine-invariant metric, so projection onto this nonempty closed
geodesically convex subset exists and is unique.

Finally, congruence by \(H^{-1/2}\) maps paths in
\(\Kron_{m,n}\) isometrically to paths in \(\mathcal S_{\Kron}(H)\). Therefore
the intrinsic distance from \(S_0\) to
\(\mathcal C_{K,\Kron}(H)\) equals the affine-invariant distance from \(G_0\)
to \(\mathcal R_{K,\Kron}(H)\). By \cref{prop:submanifold-distance}, this is
exactly \(D_{K,\Kron}(S_0;H)\).
\end{proof}

\begin{proof}[Proof of \cref{cor:kron-kstar}]
The objective \(\kappa(G^{-1}H)\) is invariant under positive rescaling of
\(G\). Hence a minimizing sequence in \(\Kron_{m,n}\) can be rescaled to satisfy
\(\det G=\det H\). Let
\[
  S=H^{-1/2}GH^{-1/2}.
\]
On this determinant gauge, \(\det S=1\). Since the sequence has bounded
condition number, the eigenvalues of \(S\) are uniformly bounded above and
below. Thus the corresponding sequence of \(S\)'s has a convergent subsequence,
and so does the sequence \(G=H^{1/2}SH^{1/2}\). The family \(\Kron_{m,n}\) is
closed by \cref{thm:kron-ai-projection}, so the limit remains Kronecker and
attains the minimum.

The equivalence
\[
  \mathcal R_{K,\Kron}(H)\ne\emptyset
  \quad\Longleftrightarrow\quad
  K\ge K_{\Kron}^\star(H)
\]
is then immediate from the definition of \(K_{\Kron}^\star(H)\).
\end{proof}

\begin{proof}[Proof of \cref{thm:kron-nearest-versus-best}]
For
\[
  \phi_{\mathrm{osc}}(x)=\max_i x_i-\min_i x_i,
\]
the spectral subdifferential at \(\log S\) is
\[
  \left\{Q_+-Q_-:
  Q_\pm\succeq0,\ \tr Q_\pm=1,\
  \operatorname{range}(Q_+)\subseteq E_+,\
  \operatorname{range}(Q_-)\subseteq E_-\right\}.
\]
By \cref{thm:normal-response-law}, the whitened subgradient for the metric
variable has the opposite sign, \(Q_--Q_+\), and global Kronecker optimality
is equivalent to one such subgradient lying in the normal space of the
Kronecker manifold. The whitened Kronecker tangent space is
\[
  \mathcal L_{m,n}
  =\{Y\otimes I_m+I_n\otimes X:X\in\mathbb S^m,\ Y\in\mathbb S^n\}.
\]
Frobenius orthogonality to \(I_n\otimes X\) for every \(X\) is equivalent to
\(\operatorname{Tr}_B(Q_--Q_+)=0\), and orthogonality to
\(Y\otimes I_m\) for every \(Y\) is equivalent to
\(\operatorname{Tr}_A(Q_--Q_+)=0\). This proves the necessary-and-sufficient
condition. When the extreme eigenvalues are simple, the only density matrices
on the two one-dimensional eigenspaces are \(u_+u_+^\top\) and
\(u_-u_-^\top\), giving the marginal formulation.

Suppose next that \(S\) has exactly two eigenvalues. Write
\[
  \log S=a\Pi_++b\Pi_-,
  \qquad r_\pm=\rank\Pi_\pm.
\]
The nearest-point normal equation gives \(\log S\in\widehat N_P\Kron\).
Because the Kronecker family is scale closed, the whitened scale direction
\(I\) is tangent, so \(\tr\log S=r_+a+r_-b=0\). Since \(a>b\), there is a
constant \(c>0\) such that
\[
  \log S
  =c\left(\frac{\Pi_+}{r_+}-\frac{\Pi_-}{r_-}\right).
\]
Taking \(Q_+=\Pi_+/r_+\) and \(Q_-=\Pi_-/r_-\) makes
\(Q_--Q_+=-c^{-1}\log S\) normal. The first part then proves global
condition optimality.

For the strict counterexample, \(\tr A=\tr B=0\), and hence
\[
  \operatorname{Tr}_B(\log H)
  =\operatorname{Tr}_B(B\otimes A)=0,
  \qquad
  \operatorname{Tr}_A(\log H)=0.
\]
The normal equations in \cref{thm:kron-ai-projection} therefore show that
\(I_6\) is the unique nearest Kronecker point. The relative log-eigenvalues at
the identity are the products of the diagonal entries of \(B\) and \(A\), so
their maximum and minimum are \(2\) and \(-2\), giving
\(\log\kappa(H)=4\).

The proposed metric commutes with \(H\), and
\[
  \log(G^{-1}H)
  =(B-\tfrac14I_3)\otimes A.
\]
The diagonal entries of \(B-\tfrac14I_3\) are
\(7/4,-3/4,-7/4\). Multiplication by the two diagonal entries \(\pm1\) of
\(A\) gives extreme relative log-eigenvalues \(7/4\) and \(-7/4\). Their
oscillation is \(7/2<4\), proving strict noncoincidence.
\end{proof}

\begin{proof}[Proof of \cref{cor:kron-two-by-two-completeness}]
Let \(P=P_{\Kron}(H)\), set \(S=P^{-1/2}HP^{-1/2}\), and write
\(L=\log S\).  The projection normal equations give
\[
  \operatorname{Tr}_A(L)=\operatorname{Tr}_B(L)=0.
\]
Introduce the real matrices
\[
  X=\begin{pmatrix}0&1\\1&0\end{pmatrix},
  \qquad
  Z=\begin{pmatrix}1&0\\0&-1\end{pmatrix},
  \qquad
  J=\begin{pmatrix}0&-1\\1&0\end{pmatrix}.
\]
The symmetric matrices \(I,X,Z\) and the skew-symmetric matrix \(J\) give an
orthogonal tensor basis for real symmetric \(4\times4\) matrices: the symmetric
tensors are generated by the nine products of \(I,X,Z\), together with
\(J\otimes J\).  The two zero-partial-trace conditions remove every term
containing an identity factor.  Hence
\[
  L
  =
  \sum_{r,s\in\{X,Z\}}C_{rs}\,r\otimes s+cJ\otimes J
\]
for a real \(2\times2\) coefficient matrix \(C\) and a scalar \(c\).

Conjugation by \(O\in O(2)\) acts on the span of \(X,Z\) through an
orthogonal \(2\times2\) transformation, and every such transformation is
realized by a suitable \(O\): planar rotations act by twice their angle on
\(\operatorname{span}\{X,Z\}\), and a reflection supplies the other connected
component of \(O(2)\).  Applying the singular value decomposition of
\(C\), choose local orthogonal matrices \(O_A,O_B\) so that conjugation by
\(O_B\otimes O_A\) transforms \(L\) into
\[
  L'=aX\otimes X+bZ\otimes Z+c'J\otimes J.
\]
The three displayed tensor products commute.  Their common orthonormal
eigenbasis is the real Bell basis
\[
\begin{aligned}
  \psi_1&=(e_1\otimes e_1+e_2\otimes e_2)/\sqrt2,&
  \psi_2&=(e_1\otimes e_1-e_2\otimes e_2)/\sqrt2,\\
  \psi_3&=(e_1\otimes e_2+e_2\otimes e_1)/\sqrt2,&
  \psi_4&=(e_1\otimes e_2-e_2\otimes e_1)/\sqrt2.
\end{aligned}
\]
Every Bell projector has both marginal density matrices equal to \(I_2/2\).
If \(\Pi_+\) and \(\Pi_-\) are the extreme spectral projectors of \(L'\),
then
\[
  Q_+=\frac{\Pi_+}{\rank\Pi_+},
  \qquad
  Q_-=\frac{\Pi_-}{\rank\Pi_-}
\]
therefore have identical marginals on both factors.  Local orthogonal
conjugation preserves this equality, so the same is true for the extreme
eigenspaces of \(L\).  \Cref{thm:kron-nearest-versus-best} proves that \(P\)
is condition-optimal.

If one factor has dimension one, the Kronecker family is the full SPD cone in
the other factor and separation is impossible.  The only remaining nontrivial
factor pair below ambient dimension six is \(2\times2\), just handled.  The
existing \(2\times3\) example attains dimension six, proving minimality.
\end{proof}

\begin{proof}[Proof of \cref{cor:kron-marginal-mismatch}]
The whitened metric-variable subdifferential of
\(G\mapsto\log\kappa(G^{-1}H)\) at \(P\) is
\[
  \mathcal W(P)
  =
  \{Q_--Q_+:Q_\pm\in\mathcal D_\pm\}.
\]
Restricting the objective to the Kronecker manifold projects this set onto the
whitened tangent space \(\mathcal L_{m,n}\), giving exactly
\(\mathcal G(P)=\Pi_{\mathcal L}\mathcal W(P)\).  The sets
\(\mathcal D_\pm\) are compact and convex, so \(\mathcal G(P)\) is compact
and convex.  Strict convexity of the squared Frobenius norm gives its unique
minimum-norm element \(g_\star\).

By the constrained Fermat rule, \(P\) is condition-optimal if and only if
\(0\in\mathcal G(P)\), equivalently \(\mu_{\mathrm{marg}}(P)=0\).  Suppose
now that \(g_\star\ne0\).  The variational characterization of the Euclidean
projection of zero onto \(\mathcal G(P)\) gives
\[
  \ip{g-g_\star}{g_\star}\ge0
  \qquad\text{for every }g\in\mathcal G(P).
\]
Hence
\[
  \min_{g\in\mathcal G(P)}\ip{g}{g_\star}
  =\fro{g_\star}^2.
\]
The one-sided directional derivative of a convex function is the support
function of its subdifferential.  Along the whitened tangent direction
\(-g_\star\), it is therefore
\[
  \max_{g\in\mathcal G(P)}\ip{g}{-g_\star}
  =-\fro{g_\star}^2.
\]
Because \(g_\star\in\mathcal L_{m,n}\), the corresponding AIRM geodesic
remains in the Kronecker manifold, proving the descent identity.

For simple extremes, both \(\mathcal D_\pm\) are singletons, which gives the
displayed formula.  With multiplicity, minimizing the norm of a linear image
of \((Q_+,Q_-)\) subject to positive-semidefinite, trace, and support
constraints is a finite-dimensional convex conic problem.
\end{proof}

\begin{proof}[Proof of \cref{thm:kron-residual-threshold}]
Let
\[
  \delta_\star=d_{\AI}(H,\Kron_{m,n})=d_{\AI}(H,P).
\]
Since \(G\in\Kron_{m,n}\), \(\delta_\star\le\rho\). By
\cref{cor:kron-residual-certificate},
\[
  0
  \le
  \frac12\rho^2-\frac12\delta_\star^2
  \le
  \frac12\eta^2,
\]
and hence \(\delta_\star\ge\sqrt{\max\{\rho^2-\eta^2,0\}}=\delta_-\). This
proves the distance bracket.

The family \(\Kron_{m,n}\) is closed under positive rescaling, so
\cref{thm:scale-closed-threshold} gives
\[
  K_{\Kron}^\star(H)
  \ge
  \exp(\alpha_{mn}\delta_\star)
  \ge
  \exp(\alpha_{mn}\delta_-).
\]
The upper bound \(K_{\Kron}^\star(H)\le\kappa_G\) holds because \(G\) is an
admissible Kronecker metric. If \(\kappa_G\le K\), then
\cref{prop:kron-relative-candidate} applied with \(G_c=G\) gives the displayed
finite \(D_{K,\Kron}\) upper bound.

It remains to prove the exact-projection condition bracket. Again by
\cref{cor:kron-residual-certificate},
\[
  d_{\AI}(G,P)\le\eta.
\]
Therefore every log-eigenvalue of \(G^{-1/2}PG^{-1/2}\) has absolute value at
most \(\eta\), and
\[
  e^{-\eta}G\preceq P\preceq e^\eta G.
\]
For any nonzero \(x\),
\[
  e^{-\eta}\frac{x^\top Hx}{x^\top Gx}
  \le
  \frac{x^\top Hx}{x^\top Px}
  \le
  e^\eta\frac{x^\top Hx}{x^\top Gx}.
\]
Taking suprema and infima over generalized Rayleigh quotients gives
\[
  \lambda_{\max}(P^{-1}H)\le e^\eta\lambda_{\max}(G^{-1}H),\qquad
  \lambda_{\min}(P^{-1}H)\ge e^{-\eta}\lambda_{\min}(G^{-1}H),
\]
so \(\kappa_P\le e^{2\eta}\kappa_G\). Interchanging \(G\) and \(P\) gives
\(\kappa_G\le e^{2\eta}\kappa_P\), equivalently
\(\kappa_P\ge e^{-2\eta}\kappa_G\).

Since \(P\) is an admissible Kronecker metric,
\(K_{\Kron}^\star(H)\le\kappa_P\).  The scale-closed lower bound and the
distance bracket give
\[
  \log K_{\Kron}^\star(H)
  \ge\alpha_{mn}\delta_\star
  \ge\alpha_{mn}\delta_-.
\]
Combining this with \(\log\kappa_P\le\log\kappa_G+2\eta\) proves the stated
projection-suboptimality chain.  The projection feasibility and
projection-infeasibility statements follow from the two-sided condition
bracket. The global impossibility certificate is the contrapositive of the
threshold lower bound.
\end{proof}

\begin{proof}[Proof of \cref{prop:kron-fixed-basis-relative}]
For
\[
  G=(V\diag(e^{b_j})V^\top)\otimes(U\diag(e^{a_i})U^\top),
\]
the matrices \(H\) and \(G\) commute in the basis \(V\otimes U\), and the
eigenvalues of \(G^{-1}H\) are
\[
  \exp(\ell_{ij}-a_i-b_j).
\]
Therefore
\[
  \log\kappa(G^{-1}H)
  =
  \max_{i,j}(\ell_{ij}-a_i-b_j)
  -
  \min_{i,j}(\ell_{ij}-a_i-b_j).
\]
The condition \(\kappa(G^{-1}H)\le K\) is thus equivalent to the existence of
a scalar \(c\) such that all residual log-eigenvalues lie in
\([c,c+\log K]\), which is exactly the displayed linear feasibility problem.
Minimizing the residual width over \(a,b\) gives the fixed-basis threshold
formula. Since \(\Kron_{U,V}\subset\Kron_{m,n}\), feasibility in the subfamily
is a full-family primal certificate; subfamily infeasibility gives no
full-family obstruction.
\end{proof}

\begin{proof}[Proof of \cref{thm:kron-fixed-basis-dual}]
For any \(Z\in\R^{m\times n}\),
\[
  \max_{ij}Z_{ij}-\min_{ij}Z_{ij}
  =
  \max_{p,q\in\Delta_{mn}}\ip{p-q}{Z},
\]
where \(\Delta_{mn}\) is the probability simplex over entries. The set
\(\{p-q:p,q\in\Delta_{mn}\}\) is exactly
\[
  \mathcal C=\{W:\sum_{ij}W_{ij}=0,\ \sum_{ij}|W_{ij}|\le2\}.
\]
Let
\[
  \mathcal A=\{A\in\R^{m\times n}:A_{ij}=a_i+b_j\}.
\]
Then
\[
  \tau^\star_{U,V}(H)=\min_{A\in\mathcal A}\max_{W\in\mathcal C}
  \ip{W}{L-A}.
\]
This is a finite-dimensional linear program. Its dual is
\[
  \max_{W\in\mathcal C\cap\mathcal A^\perp}\ip{W}{L},
\]
with no duality gap. The orthogonality condition is computed from
\[
  \ip{W}{A}
  =
  \sum_i a_i\sum_j W_{ij}
  +
  \sum_j b_j\sum_i W_{ij}.
\]
Thus \(W\in\mathcal A^\perp\) exactly when every row sum and every column sum
vanishes. These constraints imply \(\sum_{ij}W_{ij}=0\), so the dual feasible
set is the one displayed in the theorem.

If a feasible \(W\) has \(\ip{W}{L}>\log K\), then
\(\tau^\star_{U,V}(H)>\log K\), so the fixed-basis log-width target is
infeasible by \cref{prop:kron-fixed-basis-relative}. For any nonzero
zero-total \(W\), the identity
\(\sum W^+=\sum W^-=\frac12\norm{W}_1\) gives the normalized witness form when
\(\norm{W}_1=2\).
\end{proof}

\begin{proof}[Proof of \cref{thm:kron-certify-soundness}]
\Cref{prop:kron-fixed-basis-relative} proves soundness of the fixed-basis
feasible status.  \Cref{thm:kron-fixed-basis-dual} proves the obstruction
status.  The four full noncommuting statuses follow from
\cref{thm:kron-residual-threshold}. If none of those conditions holds, the
procedure records only \(\mathrm{INCONCLUSIVE}\), which makes no claim.
\end{proof}

\begin{proof}[Proof of \cref{cor:kron-certify-interval}]
The enclosure assumptions imply
\[
  \delta_-
  =
  \sqrt{\max\{\rho^2-\eta^2,0\}}
  \ge
  \sqrt{\max\{\underline\rho^2-\overline\eta^2,0\}}
  =\underline\delta.
\]
Therefore
\(K_{\Kron}^\star(H)\ge\exp(\alpha_{mn}\underline\delta)\), proving the
first implication.  The second follows from
\(\kappa_G\le\overline\kappa\).  For the exact projection, the bracket in
\cref{thm:kron-residual-threshold} and the enclosures give
\[
  e^{-2\overline\eta}\underline\kappa
  \le
  e^{-2\eta}\kappa_G
  \le
  \kappa_P
  \le
  e^{2\eta}\kappa_G
  \le
  e^{2\overline\eta}\overline\kappa.
\]
The remaining two implications follow immediately.  Returning
\(\mathrm{INCONCLUSIVE}\) when no strict certified comparison holds makes no
claim and is therefore sound.
\end{proof}

\begin{proof}[Proof of \cref{prop:kron-relative-candidate}]
The matrix
\[
  S_c^{-1}=H^{1/2}G_c^{-1}H^{1/2}
\]
has the same eigenvalues as \(G_c^{-1}H\), so
\(\kappa(S_c)=\kappa(G_c^{-1}H)\). Hence the displayed condition is exactly
membership of \(S_c\) in \(\Ck\), and \(S_c\in\mathcal C_{K,\Kron}(H)\).

By \cref{thm:kron-ai-projection}, \(\Kron_{m,n}\) is totally geodesic in the
ambient affine-invariant SPD cone. Therefore the affine-invariant geodesic
from \(G_0\) to \(G_c\) remains in \(\Kron_{m,n}\) and has length
\(d_{\AI}(G_0,G_c)\). Congruence by \(H^{-1/2}\) is an isometry for the
affine-invariant metric, so the relative path
\[
  S_t=H^{-1/2}G_tH^{-1/2}
\]
is an admissible path in \(\mathcal S_{\Kron}(H)\) from \(S_0\) to the target
point \(S_c\) with the same length. Taking the infimum over all admissible
paths gives
\[
  D_{K,\Kron}(S_0;H)\le d_{\AI}(G_0,G_c).
\]
The closed-form expression for the right-hand side is
\cref{prop:kron-product-distance}. The last statement is only the logical
one-way nature of a sufficient certificate.
\end{proof}

\begin{proof}[Proof of \cref{prop:kron-kkt}]
The objective gradient at \(Q\), in whitened coordinates and restricted to
\(\mathcal L_{m,n}\), is \(-\Pi_{\mathcal L}R_Q(M)\). For the path
\[
  Q(t)=Q^{1/2}\exp(tZ)Q^{1/2},
\]
the directional derivative of \(\log\lambda_{\max}\) is
\[
  \max_{\substack{u\in E_{\max}\\ \norm{u}=1}}u^\top Zu,
\]
and the directional derivative of \(\log\lambda_{\min}\) is
\[
  \min_{\substack{v\in E_{\min}\\ \norm{v}=1}}v^\top Zv.
\]
Thus the directional derivative of \(\omega=\log\lambda_{\max}-\log\lambda_{\min}\)
is the support function of the stated subdifferential. Since \(K>1\) has the
strict feasible point \(I_{mn}\), the convex KKT condition on the tangent
Hadamard submanifold gives
\[
  0\in-\Pi_{\mathcal L}R_\star+\mu\Pi_{\mathcal L}\partial\omega(Q_\star),
  \qquad \mu\ge0,
\]
with complementarity. This is the displayed equation; inactive constraints
have \(\mu=0\).
\end{proof}

\begin{proof}[Proof of \cref{prop:kfac-certificates}]
Projection onto the closed geodesically convex family \(\Kron_{m,n}\) satisfies
the Hadamard Pythagorean inequality. Applying it to
\[
  H_t,\qquad P_t=P_{\Kron}(H_t),\qquad G_t\in\Kron_{m,n}
\]
gives \(E_t^2\ge M_t^2+A_t^2\). Applying
\cref{thm:kron-k-target} with \(M=H_t\) and \(G_\star=P_t\) gives the two-sided
self-conditioned \(K\)-target bound. The computability statements are exactly
\cref{thm:kron-ai-projection,thm:kron-k-target,prop:kron-product-distance}.
\end{proof}

\begin{proof}[Proof of \cref{prop:proxy-inflation}]
From
\[
  (1-\varepsilon)H\preceq \widehat H\preceq (1+\varepsilon)H
\]
we obtain
\[
  \frac{1}{1+\varepsilon}
  \frac{v^\top\widehat H v}{v^\top Gv}
  \le
  \frac{v^\top H v}{v^\top Gv}
  \le
  \frac{1}{1-\varepsilon}
  \frac{v^\top\widehat H v}{v^\top Gv}.
\]
Taking minima and maxima over nonzero \(v\) gives the bound on generalized
condition numbers.
\end{proof}

\begin{proof}[Proof of \cref{prop:stoch-proxy-upper}]
On the stated Loewner event, \cref{prop:proxy-inflation} gives
\[
  \kappa(G^{-1}H)
  \le
  \frac{1+\varepsilon}{1-\varepsilon}
  \kappa(G^{-1}\widehat H)
  \le
  \frac{1+\varepsilon}{1-\varepsilon}\widehat K
  =K.
\]
The event has probability at least \(1-\delta\), which proves the claim.
\end{proof}
